\appendixitem{实代数的表示}

记 \(\sigma\) 为 C 的自同构 \(\alpha\mapsto\bar{\alpha}\)；若 W 是复向量空间，则记 \(\overline{W}\) 为 C-向量空间 \(\sigma_{*}(W)\)（即，W 这个群赋予作用律 \((\alpha,w)\mapsto\bar{\alpha}w\)，其中 \(\alpha\in C, w\in W\)）。

命题 1. 设 A 是一个 R-代数（结合且有单位），V 是 R 上有限维单 A-模。则必属下列三种情形之一：

\(\alpha)\) V 的交换子（Algebra, Chap. VIII, § 5, no. 1）同构于 \(\mathbf{R}\)，且 \(\mathrm{A}_{(\mathbf{C})}\)-模 \(\mathrm{V}_{(\mathbf{C})}\) 是单模；

\(\beta)\) V 的交换子同构于 \(\mathbf{C}\)；\(\mathrm{A}_{(\mathbf{C})}\)-模 \(\mathrm{V}_{(\mathbf{C})}\) 是两个不相同构的单 \(\mathrm{A}_{(\mathbf{C})}\)-模的直和，而 \(\sigma \otimes 1_{\mathrm{V}}\) 交换这两个子模；

\(\gamma)\) V 的交换子同构于 H；\(A_{(C)}\)-模 \(V_{(C)}\) 是两个同构的单 \(A_{(C)}\)-模的直和，而 \(\sigma \otimes 1_V\) 交换这两个子模。

V 的交换子 E 是一个域，是 R 的有限扩域（Algebra, Chap. VIII, §3, no. 2, Prop. 2），因而同构于 R、C 或 H（Algebra, Chap. VIII, §15）。\(A_{(C)}\)-模 \(V_{(C)}\) 是半单的（Algebra, Chap. VIII, §11, no. 4），其交换子可与 \(C \otimes_{R} E\) 识别（Algebra, Chap. VIII, §11, no. 2, Lemma 1）。

若 \(\mathbf{E}\) 同构于 \(\mathbf{R}\)，则 \(V_{(\mathbf{C})}\) 的交换子同构于 \(\mathbf{C}\)，且 \(V_{(\mathbf{C})}\) 是单 \(A_{(\mathbf{C})}\)-模（Algebra, Chap. VIII, §11, no. 4）。

若 \(\mathbf{E}\) 不同构于 \(\mathbf{R}\)，则它包含一个同构于 \(\mathbf{C}\) 的域；于是 \(V\) 具有一个 \(A_{(\mathbf{C})}\)-模结构，记为 \(V^c\)。则 \(V^c\) 是一个单

\(A_{(C)}\)-模，且 C-线性映射 \(\psi: V_{(C)} \to V^{c} \oplus \overline{V}^{c}\) 满足 \(\psi(\alpha \otimes v) = (\alpha v, \bar{\alpha} v)\)，其中 \(\alpha \in C\)、\(v \in V\)；该映射是同构（Algebra, Chap. V, §10, no. 4, Prop. 8）。此外，在此同构下，\(\sigma \otimes 1_{V}\) 对应于 R-自同构 \((v, v') \mapsto (v', v)\)（作用在 \(V^{c} \oplus \overline{V}^{c}\) 上），因而交换两个 \(A_{(C)}\)-子模 \(\psi^{-1}(V^{c})\) 和 \(\psi^{-1}(\overline{V}^{c})\)。

因此，\(E_{(\mathbf{C})}\)（\(V_{(\mathbf{C})}\) 的交换子）包含 \(C \times C\)，后者通过数乘变换作用于 \(V^{c} \oplus \overline{V}^{c}\)。不存在一个 \(A_{(\mathbf{C})}\)-模同构，将 \(V^{c}\) 映到 \(\overline{V}^{c}\)，当且仅当 \(E_{(\mathbf{C})}\) 退化为 \(C \times C\)，也就是当且仅当 E 同构于 C。证明完毕。

命题 2. 设 A 是一个 R-代数（结合且有单位），W 是 C 上有限维单 \(\mathrm{A}_{(\mathbf{C})}\)-模。则必属下列三种情形之一：

\begin{enumerate}
\def\labelenumi{\alph{enumi})}
\item
  存在一个 \(A_{(\mathbf{C})}\)-同构 \(\theta\) 从 W 到 \(\overline{W}\)，满足 \(\theta \circ \theta = 1_{W}\)。则 \(\theta\) 的不动点集 V 是 W 上的一个 R-结构，并且是交换子为 R.1 \(_{V}\) 的单 A-模。此外，\(W_{[R]}\) 是两个同构单 A-模的直和。
\item
  \(\mathrm{A}_{(\mathbf{C})}\)-模 W 与 \(\overline{\mathrm{W}}\) 不同构；则 \(\mathrm{W}_{[\mathbf{R}]}\) 是交换子为 \(\mathbf{C}.1_{\mathrm{W}}\) 的单 A-模。
\item
  存在一个 \(\mathrm{A}_{(\mathbf{C})}\)-同构 \(\theta\) 从 W 到 \(\overline{W}\)，满足 \(\theta \circ \theta = -1_{W}\)。则 A-模 \(W_{[R]}\) 是单模，其交换子是域 \(C.1_{W} \oplus C.\theta\)，同构于 H。
\end{enumerate}

复向量空间 \(\mathrm{Hom}_{\mathrm{A}_{(\mathbf{C})}}(\mathrm{W},\overline{\mathrm{W}})\) 的维数 \(\leq 1\)（Algebra, Chap. VIII, §3, no. 2）；若 \(\theta \in \mathrm{Hom}_{\mathrm{A}_{(\mathbf{C})}}(\mathrm{W},\overline{\mathrm{W}})\)，则 W 的自同态 \(\theta \circ \theta\) 是数乘变换，其倍数为 \(\alpha \in \mathbf{C}\)。对于所有 \(w \in \mathrm{W}\)，有 \(\alpha \theta(w) = \theta \circ \theta \circ \theta(w) = \theta(\alpha w) = \bar{\alpha} \theta(w)\)，所以 \(\alpha\) 是实数。若 \(\theta' = \lambda \theta\)，其中 \(\lambda \in \mathbf{C}\)，则 \(\theta' \circ \theta' = |\lambda|^2 \theta \circ \theta\)；因此，下列三种可能性中恰有一种成立：

\begin{enumerate}
\def\labelenumi{\alph{enumi})}
\item
  存在 \(\theta \in \mathrm{Hom}_{\mathrm{A}_{(\mathbf{C})}}(\mathrm{W},\overline{\overline{\mathrm{W}}})\)，满足 \(\theta \circ \theta = 1_{\mathrm{W}}\)；
\item
  \(\mathrm{Hom}_{\mathrm{A}_{(\mathbf{C})}}(\mathrm{W},\overline{\mathrm{W}}) = \{0\}\)；
\item
  存在 \(\theta \in \mathrm{Hom}_{\mathrm{A}_{(\mathbf{C})}}(\mathrm{W},\overline{\mathrm{W}})\)，满足 \(\theta \circ \theta = -1_{\mathrm{W}}\)。
\end{enumerate}

在情形 a) 中，\(\theta\) 的不动点集 V 是 W 上的一个 R-结构（Algebra, Chap. V, p.~61, Prop. 7）；由于 \(V_{(\mathbf{C})}\) 同构于 W，A-模 V 是交换子为 \(R.1_{V}\) 的单模（命题 1），而 \(W_{[R]}\) 不是单模。

反之，若 \(W_{[R]}\) 不是单模，设 V 是 \(W_{[R]}\) 的一个单 A-子模；由于 \(A_{(C)}\)-模 W 是单模，有 \(V + iV = W\) 且 \(V \cap iV = \{0\}\)，即 \(W = V \oplus iV\)。因此，V 是 W 上的一个 R-结构，而同构 \(\theta\) 从 W 到 \(\overline{W}\) 满足：有 \(\theta(v + iv') = v - iv'\)，其中 v 和 \(v'\) 属于 V，并且 \(\theta \circ \theta = 1_{W}\)。

因此，在情形 \(b)\) 和 \(c)\) 中，A-模 \(\mathrm{W}_{[\mathbf{R}]}\) 是单模；由命题 1，其交换子 E 在情形 \(b)\) 中同构于 C，在情形 \(c)\) 中同构于 H。此外，显然 E 包含 \(\mathbf{C}.1_{\mathrm{W}}\)，并包含 \(\mathbf{C}.\theta\)（在情形 \(c)\) 中），命题得证。

在命题的假设下，当情形分别为 a)、b) 或 c) 时，称 \(A_{(C)}\)-模 W（相对于 A）分别为实型、复型或四元数型。

对于 K = R 或 C，记 \(\mathfrak{S}_{\mathrm{K}}(\mathrm{A})\) 为 K 上有限维单 \(A_{(K)}\)-模的等价类集合。群 \(\Gamma = \operatorname{Gal}(\mathbf{C}/\mathbf{R})\) 作用于 \(\mathfrak{S}_{\mathbf{C}}(\mathrm{A})\)；前面两个命题在 \(\mathfrak{S}_{\mathbf{R}}(\mathrm{A})\) 与商集 \(\mathfrak{S}_{\mathbf{C}}(\mathrm{A})/\Gamma\) 之间建立了一个双射对应。
% semantic-fragments: legacy-input-seam
\relax{}
